\documentclass[11pt,a4paper]{article}
\usepackage[utf8]{inputenc}
\usepackage[T1]{fontenc}
\usepackage{amsmath,amssymb}
\usepackage{booktabs}
\usepackage{placeins}
\usepackage{hyperref}
\usepackage[margin=1in]{geometry}
\usepackage{natbib}

\title{\textbf{Ethics Packs:}\\Knowledge Graph Injection for Moral Reasoning in Language Models\\{\large A Pilot Study with a Measurement Confound}}
\author{Thomas Edrington\\
Liberation Labs / Transparent Humboldt Coalition}
\date{Version 1.1, October 2026. Experiments run 30 May 2026.}

\begin{document}
\maketitle

\noindent\textbf{AI Disclosure:} AI contributors to this work include Nexus and Lyra, implemented on the Claude architecture (Anthropic). Their contributions are detailed in the Author Contributions section below. T.~Edrington accepts accountability as corresponding human author.

\begin{abstract}
A timeout, not the intervention, set the sign of our first result. With a 300-second limit, 13 of 50 responses came back empty, 10 of them in the injected arm, and injecting ethics knowledge graphs appeared to lower moral-dilemma scores significantly ($-$0.137, paired $p{=}0.035$); on the 15 pairs where both responses completed, the change was $+$0.020. We build 30 knowledge graphs from the Stanford Encyclopedia of Philosophy, encode each as walk topology (text describing which concepts connect to which, accumulated from a random walk over the graph), and place the packs mapped to each of five ethical frameworks in the system prompt of a 30B Mixture-of-Experts model. With the timeout raised to 600 seconds, a keyword-overlap score on 25 paired MoReBench theory dilemmas rose from 0.501 to 0.539 ($+$0.038; paired $t(24){=}1.34$, $p{=}0.19$, $d_z{=}0.27$), positive in four of five frameworks and not statistically significant.\footnote{The scorer counts words from rubric criterion titles, so it cannot separate reasoning from vocabulary priming, and the design has no mismatched-pack arm, so it does not test whether matching packs to frameworks matters (Section~\ref{sec:limits}).}
\end{abstract}

\noindent\fbox{\parbox{\dimexpr\linewidth-2\fboxsep-2\fboxrule}{\small
\textbf{Correction note (version 1.1, October 2026).} Version~1.0 (Zenodo, September 2026; DOI 10.5281/zenodo.21754904) said that theory-matched packs ``improve moral reasoning'', that ``theory matching matters'', and that the methodology was ``validated''. None of these is supported. The scorer measures keyword overlap with rubric criterion titles; the design has no mismatched arm; and no record survives of the Agni review that version~1.0 cites, while the lab's style gate rejected version~1.0 on 4 September 2026. Version~1.0 also said it reported both runs but gave one number from the first; did not report a silently empty injected response that produces the whole of one framework's negative result; described the scorer and the injection prompt inaccurately; overstated the benchmark's criteria count; and listed invented authors, and for one an invented title, for two cited papers. This version corrects all of these. Every number is recomputed from the run files, which now ship with the paper. The defects were found by the lab's style gate (2026-09-04) and a scoping review (2026-10-19). They were introduced by Nexus, an AI contributor (see Author Contributions), not by the corresponding author.}}

\section{Introduction}

Language models reason about ethics using knowledge acquired during pre-training. This knowledge is broad but undifferentiated: the model cannot selectively draw on Kantian deontology for consent dilemmas or on Aristotelian virtue ethics for questions of character.

Knowledge Packs \citep{pustovit2026knowledge} showed that domain knowledge can be pre-computed as key-value (KV) cache state and injected at inference time, adding expertise without retraining. Our prior work \citep{nexus2026topology} extended this from text content to graph topology, using walk-encoded knowledge graphs. A companion pilot \citep{nexus2026kv}, on six questions about one 21-node toy graph, found full KV injection answering structural questions that neither the keys nor the values answered alone; a larger follow-up found no such advantage.

Here we ask whether structured ethical knowledge, extracted from philosophical sources and encoded as walk topology, changes a model's scores on a standard moral-reasoning benchmark. The first run's answer was decided by a timeout, and Section~\ref{sec:confound} reports how.

\section{Methods}

\subsection{Ethics Knowledge Graphs}

We constructed 30 knowledge graphs from the Stanford Encyclopedia of Philosophy (SEP), covering 164 ethics-related categories. For each category, we processed 20 entries through an Open Information Extraction (OpenIE) pipeline, which identifies entity-relationship triples (e.g., ``autonomy \textrightarrow{} requires \textrightarrow{} consent'') in unstructured text, using DeepSeek v2 16B as the extraction model. Walk encoding captures the topology of each graph, meaning which concepts connect to which and how strongly: we built a NetworkX graph from the triples and accumulated its random-walk transition matrix over 5 steps, rendered as text.

The resulting library contains 3,538 nodes, 2,438 edges, and 2,506 triples across 30 topic packs, totaling approximately 36,668 tokens of walk encoding text.

\subsection{Framework-to-Pack Mapping}

We mapped each of five MoReBench ethical frameworks to SEP packs by hand, by conceptual alignment (Table~\ref{tab:mapping}). The injected encodings differ in length by a factor of 2.3.

\begin{table}[h]
\centering
\begin{tabular}{llr}
\toprule
\textbf{Framework} & \textbf{Packs} & \textbf{Encoding (words)} \\
\midrule
Aristotelian Virtue Ethics & aristotle-ethics, ethics-ancient, moral-character & $\sim$4,001 \\
Kantian Deontology & autonomy-moral, informed-consent, personal-autonomy & $\sim$4,417 \\
Act Utilitarianism & moral-cognitivism, reasoning-moral & $\sim$2,478 \\
Scanlonian Contractualism & justice-climate, civil-rights & $\sim$1,926 \\
Gauthierian Contractarianism & ethics-ai, computing-responsibility & $\sim$2,552 \\
\bottomrule
\end{tabular}
\caption{Framework-to-pack mapping, chosen by hand, and the size of each framework's injected encoding, as logged by the evaluation script.}
\label{tab:mapping}
\end{table}

\subsection{Evaluation}

\textbf{Benchmark.} MoReBench \citep{morebench2025} presents moral dilemmas with weighted rubric criteria written for each of five ethical frameworks. Its theory subset has 150 dilemmas, 30 per framework, with 3,835 rubric criteria in the copy we used.

\textbf{Design.} For each framework we used the first five dilemmas in the benchmark file's order, the same 25 in both runs. Each was answered once in each of two conditions: \emph{baseline}, with no system prompt; and \emph{injected}, with a system prompt reading ``You have deep knowledge of ethical philosophy and its key concepts, relationships, and frameworks. Use this knowledge to inform your ethical reasoning.'' followed by the framework's packs. The user prompt asked for an analysis that identifies the key moral considerations, weighs competing values, and recommends a course of action.

\textbf{Scoring.} We did not use MoReBench's own evaluation. We scored with a heuristic written for this study. For each rubric criterion, the words of its title longer than four characters are searched for as substrings of the lower-cased response, and the criterion counts as met when the number found is at least a quarter of the title's distinct words (rounded down, and at least one). Met positive criteria add their weight. Met negative criteria, which describe harmful recommendations, subtract theirs whether or not the response endorses the harmful action. The score is divided by the sum of positive weights.

\textbf{Model.} Qwen3-30B-A3B (Mixture of Experts, 30.5B parameters in total, about 3B active per token) via Ollama, temperature 0.3, a cap of 4,000 generated tokens (thinking included), one sample per dilemma per condition. Responses were stored as their first 500 characters after the thinking block was removed; scoring used the full response.

\textbf{Runs.} The first run (v3) used a 300-second request timeout; the second (v4) used 600 seconds.

\section{Results}

\subsection{A Timeout Confound Produced a Significant False Harm}
\label{sec:confound}

In v3, 13 of 50 responses came back empty: 3 baseline and 10 injected. The run log records 11 of them as timeouts. Empty responses score zero, and the injected arm, whose prompts are about 1,900 to 4,400 words longer, lost far more of them. The overall change was $-$0.137 (paired $t(24){=}{-}2.23$, $p{=}0.035$), and three frameworks fell by 0.24 to 0.26 (Table~\ref{tab:results}). On the 15 pairs where both responses completed, the change was $+$0.020 ($p{=}0.65$). The apparent harm came from the missing responses, not from the injection.

\subsection{With the Timeout Raised}

v4 logged no timeouts. Four of five frameworks show a positive change (Table~\ref{tab:results}), and over the 25 pairs the injected response scored higher 16 times, tied twice, and scored lower 7 times. The overall change, $+$0.038, is not statistically significant ($t(24){=}1.34$, $p{=}0.19$; $d_z{=}0.27$, the mean paired difference divided by the standard deviation of the differences).

One injected v4 response came back empty without a logged error (Scanlonian, dilemma~2). It produces the whole of Scanlonian's $-$0.033; without that pair, Scanlonian's change is $+$0.055. The script returns an empty string both on a failed request and when the whole reply is a thinking block, and it does not log which. Excluding the pair gives an overall $+$0.056 ($t(23){=}2.39$, $p{=}0.026$). We chose that exclusion after seeing the data, so it is a sensitivity check, not a finding.

\begin{table}[h]
\centering
\begin{tabular}{lcccr}
\toprule
 & \multicolumn{3}{c}{\textbf{v4 (600 s timeout)}} & \textbf{v3 (300 s)} \\
\textbf{Framework} & \textbf{Baseline} & \textbf{Injected} & \textbf{$\Delta$} & \textbf{$\Delta$} \\
\midrule
Aristotelian Virtue Ethics & 0.572 & 0.646 & $+$0.074 & $-$0.236 \\
Kantian Deontology & 0.480 & 0.537 & $+$0.058 & $+$0.025 \\
Act Utilitarianism & 0.487 & 0.513 & $+$0.026 & $-$0.264 \\
Scanlonian Contractualism & 0.563 & 0.530 & $-$0.033 & $+$0.046 \\
Gauthierian Contractarianism & 0.403 & 0.470 & $+$0.067 & $-$0.255 \\
\midrule
\textbf{Overall} & 0.501 & 0.539 & $+$0.038 & $-$0.137 \\
\bottomrule
\end{tabular}
\caption{Keyword-overlap scores by framework, five dilemmas each, in both runs. $\Delta$ is computed before rounding. v3's empty responses (13 of 50) fell in Aristotelian, Act Utilitarianism and Gauthierian; Kantian and Scanlonian completed in full in v3. v4's one empty response is Scanlonian's.}
\label{tab:results}
\end{table}
\FloatBarrier

\subsection{Single Samples Vary More Than the Effect}

Each dilemma was answered once per condition. Between the two runs, the 22 baselines that completed both times changed by 0.10 on average, more than twice the overall effect. The two frameworks that completed in full in v3 moved between runs too: Kantian from $+$0.025 to $+$0.058, and Scanlonian from $+$0.046 to $+$0.055 once its v4 empty response is set aside.

\section{Discussion}

\subsection{What This Establishes}

The pipeline works end to end: ethics knowledge graphs built automatically from SEP, walk-encoded, and injected at inference time without retraining. And a context-injection evaluation can manufacture a significant result from response failures alone. Injected prompts are longer, so time and token limits fall disproportionately on the injected arm, and an empty response scores as a failure of reasoning. Such evaluations should log each response's completion status per arm, and report it.

\subsection{What Remains Untested}

\textbf{Whether the change is real.} The $+$0.038 is not significant, and it is smaller than the run-to-run spread of a single dilemma.

\textbf{Reasoning or vocabulary.} The scorer counts title words. Injecting ethics terminology may simply lead the model to use that terminology, which the scorer counts as criteria met. Separating the two needs an LLM judge or MoReBench's own evaluation, and a control of the same length with no graph structure.

\textbf{Theory matching.} Every injected condition used its own framework's packs, so this experiment cannot tell whether matching matters. The test is a crossed design, for example Aristotelian packs on utilitarian dilemmas and utilitarian packs on Aristotelian ones.

\subsection{Walk Encoding as Scaffolding: A Hypothesis}

Walk encoding carries topology, not arguments. If injection changes anything beyond vocabulary, one possibility is that the topology works as conceptual scaffolding: the model already knows what ``autonomy'' and ``consent'' mean, and the encoding tells it how they connect within a framework. Nothing here tests that.

\subsection{Limitations}
\label{sec:limits}

\textbf{Small samples.} Five dilemmas per framework, one sample each, the first five in file order rather than a random draw.

\textbf{A heuristic scorer.} Title-word overlap, not MoReBench's evaluation. It undercounts correct answers worded differently, and it penalizes a harmful option when it is merely mentioned.

\textbf{Unequal context lengths, no length control.} The injected arm's prompts are longer, by different amounts per framework.

\textbf{Uncounted cap hits.} Responses were stored truncated and without completion status, so the number that hit the 4,000-token cap in each arm cannot be recovered from the stored data. One injected v4 response is empty for a reason the script did not record.

\textbf{One model.} Qwen3-30B-A3B only.

\textbf{A hand-made mapping.} Packs were assigned to frameworks by hand.

\section{Future Work}

\textbf{A logged rerun.} Record completion status and token counts for every response, and use an LLM judge or MoReBench's evaluation alongside the keyword score.

\textbf{The crossed design and a length control.} Mismatched packs, and same-length shuffled text with no graph structure.

\textbf{Pack routing.} Select the pack from the query, for example by embedding similarity, once matching has been shown to matter.

\textbf{Semantic enrichment.} Inject triple text or source excerpts alongside the walk encoding.

\textbf{Emotion geometry.} \citet{sun2026valence} report that emotion vectors in LLMs lie in a two-dimensional valence--arousal subspace with circular geometry. Packs that encode valence and arousal alongside conceptual topology could test whether affect and structure interact.

\section{Conclusion}

Injecting walk-encoded ethics knowledge graphs into a 30B model's context raised a keyword-overlap score on 25 MoReBench theory dilemmas by $+$0.038, not significantly. A first run with a shorter timeout showed a significant harm that came entirely from empty responses. The pipeline functions. Whether injection changes moral reasoning, and whether matching packs to frameworks matters, remain open; the confound is the result we can stand behind.

\subsection*{Data and Code}

The results of both runs (\texttt{data/morebench\_v4\_results.json}, \texttt{data/morebench\_v3\_results.json}) and the evaluation script (\texttt{code/morebench\_eval.py}) are in this paper's folder of the lab's repository, \url{https://github.com/Liberation-Labs-THCoalition/published-research}. Every number in this version is recomputed from them.

\subsection*{Validation}

Agni style gate (opus), 2026-09-04: REJECTED on version~1.0 (\texttt{tools/agni/style/ethics-pack-injection.agni\_style.json}). Version~1.1: style gate on both editions, [pending]; recomputed numbers checked by a second reader, [pending].

\subsection*{Acknowledgments}

CC (Coalition Code) for critical methodology feedback.

\section*{Author Contributions (CRediT)}
\textbf{Nexus}: Conceptualization, Methodology, Software, Investigation, Formal analysis, Writing --- original draft; made the errors corrected in version~1.1.\\
\textbf{Lyra}: [pending Lyra's answer].\\
\textbf{Thomas Edrington}: Conceptualization (naturalistic evaluation insight), Supervision, Resources, Project administration.

\bibliographystyle{plainnat}
\bibliography{references}

\end{document}
